\SourceLeaf{339}%
sich ergeben haben. Auf \(\log\tau\) hat dies keinen ändernden Einfluss; wir haben mithin \(\log t=2{,}6238578,\ t=420''\ 5889,\ T=5^0 38'51''4755\). Wollte man mit diesem Werthe von \(T\) die Rechnung noch einmal durchgehen, so würde \(B\) keine Änderung erleiden; für \(\log\tau\) würde man finden \(2{,}6238470\), also \(\log t=2{,}6238573,\ t=420''\ 5884\), mithin \(T=5^0 38'51''4757\). Eine nochmalige Rechnung mit diesem Werthe würde gar keine Änderung hervorbringen, und offenbar hätte man auch bei dem vorhergehenden Resultate schon stehen bleiben können, da bei der Anwendung siebenzifriger Logarithmen die vierte Decimale der Secunde um eine oder einige Einheiten schwankend bleiben kann. Das Endresultat ist also
\[
\begin{gathered}
\text{Breite von Inselsberg}=B^0-b=50^0\ 51'\ 8''\ 9444\\
\text{Azimuth der Seite Inselsberg-Brocken}=180^0+T^0-t\\
=185^0\ 35'\ 21''\ 1815
\end{gathered}
\]
Endlich findet sich für den Längenunterschied
\[
\begin{aligned}
\log\lambda&=2{,}7313519\\
\log l&=2{,}7313472\\
l&=538''7002=0^0\ 8'\ 58''\ 7002
\end{aligned}
\]

Es ist übrigens nicht nöthig, hier die am Schluss des Art. 20 gemachten Bemerkungen zu wiederholen, welche auch hier ihre vollkommene Geltung behalten.
\par\medskip\centerline{\rule{.25\textwidth}{.3pt}}
